## 1. The Core Concept

ANOVA tests the null hypothesis that the means of $k$ populations are equal. It does this by comparing the variance **between** groups to the variance **within** groups.

* **Null Hypothesis ($H_0$):** $\mu_1 = \mu_2 = \dots = \mu_k$
* **Alternative Hypothesis ($H_a$):** At least one $\mu_i$ is different.

### Fundamental Identity

The Total Sum of Squares ($SS_{Total}$) is decomposed as:


$$SS_{Total} = SS_{Between} + SS_{Within}$$


(Also written as $SS_{Total} = SS_{Treatment} + SS_{Error}$)

---

## 2. One-Way ANOVA Table

This is the heart of any ANOVA problem. You will likely be asked to fill one of these out.

| Source | Deg. of Freedom ($df$) | Sum of Squares ($SS$) | Mean Square ($MS$) | F-Statistic |
Between | $k - 1$ | $SS_B$ | $MS_B = \frac{SS_B}{k-1}$ | $F = \frac{MS_B}{MS_W}$ |
Within | $N - k$ | $SS_W$ | $MS_W = \frac{SS_W}{N-k}$ |  |
Total | $N - 1$ | $SS_T$ |  |  |

* $N$ = Total number of observations.
* $k$ = Number of groups/levels.

---

## 3. Assumptions (The "LINE" mnemonic)

For the F-test to be valid, the residuals ($\epsilon_{ij}$) must satisfy:

1. **L**inearity/Additivity: The model is a linear combination of effects.
2. **I**ndependence: Observations are independent (crucial for the $df$ calculation).
3. **N**ormality: Residuals are normally distributed, $\epsilon_{ij} \sim N(0, \sigma^2)$.
4. **E**qual Variance (Homoscedasticity): $\sigma_1^2 = \sigma_2^2 = \dots = \sigma_k^2$.

---

## 4. Key Formulas for $SS$

If you have to calculate these by hand, use the "Computational Forms":

* **$SS_{Total} = \sum X^2 - \frac{(\sum X)^2}{N}$**
* **$SS_{Between} = \sum \frac{T_i^2}{n_i} - \frac{(\sum X)^2}{N}$** (where $T_i$ is the sum of group $i$)
* **$SS_{Within} = SS_{Total} - SS_{Between}$**

---

## 5. Post-Hoc Tests

If you reject $H_0$, ANOVA doesn't tell you *which* groups differ. You must use post-hoc tests to control the **Family-Wise Error Rate (FWER)**.

* **Tukey’s HSD:** Used for all pairwise comparisons when group sizes are equal.
* **Bonferroni Correction:** Very conservative; divide $\alpha$ by the number of tests ($m$): $\alpha_{new} = \frac{\alpha}{m}$.
* **Scheffé’s Method:** Most flexible, used for "unplanned" or complex contrasts.

---

## 6. Effect Size: $\eta^2$ (Eta-Squared)

While the p-value tells you if there is an effect, $\eta^2$ tells you the magnitude.


$$\eta^2 = \frac{SS_{Between}}{SS_{Total}}$$


*Interpret as the proportion of total variance explained by the grouping factor.*

---

## Quick Exam Tips

* **The F-distribution:** Always right-skewed. The F-statistic is always positive. If $F \approx 1$, you fail to reject $H_0$.
* **Relationship to t-test:** For $k=2$, $F = t^2$.
* **Mean Squares:** Note that $MS_W$ (or $MS_{Error}$) is an unbiased estimator of the common population variance $\sigma^2$.

Would you like me to create a similar summary for **Two-Way ANOVA** or perhaps work through a numerical example to practice the $SS$ calculations?

The \( F \)-distribution is the distribution of a random variable
\( X = \frac{\chi^2_1 / d_1}{\chi^2_2 / d_2} \), where \( \chi^2_{1}, \chi^2_{1} \) are independent chi-square distributions with df \( d_1, d_2 \) respectively.



**Factor Analysis (FA)** is best understood as a statistical method for **data reduction** and **latent variable modeling**. 
Unlike PCA (which is a linear combination of observed variables), FA assumes that observed variables are linear combinations of unobserved (latent) factors plus error.
## 1. The Fundamental Equation

In FA, we model a vector of $p$ observed variables $X$ as a linear combination of $m$ common factors $F$ and $p$ specific errors $\epsilon$.

$$X = \Lambda F + \epsilon$$

* $X_{p \times 1}$: Vector of observed variables ($X - \mu$).
* $\Lambda_{p \times m}$: Matrix of **Factor Loadings** ($\lambda_{ij}$ is the loading of variable $i$ on factor $j$).
* $F_{m \times 1}$: Vector of **Common Factors** ($E(F)=0, Var(F)=I$).
* $\epsilon_{p \times 1}$: Vector of **Specific Factors** (Unique variance/error).

### Covariance Structure

The model implies the following decomposition of the covariance matrix $\Sigma$:
$$\Sigma = \Lambda \Lambda^T + \Psi$$
Where $\Psi = diag(\psi_1, \dots, \psi_p)$ is the variance of the specific factors.

## 2. Key Terminology
* **Communality ($h_i^2$):** The proportion of variance in variable $i$ explained by the $m$ common factors.
$$h_i^2 = \sum_{j=1}^m \lambda_{ij}^2$$
* **Uniqueness ($\psi_i$):** The variance "unique" to variable $i$ (not shared with others).
$$\psi_i = 1 - h_i^2$$

* **Eigenvalues:** In the context of FA, these represent the amount of variance explained by each factor.
* **Factor Score:** The estimated value of a latent factor for a specific observation.

## 3. Exploratory Factor Analysis (EFA) Process

### Step A: Check Factorability

Before running FA, ensure your data is suitable:

* **Kaiser-Meyer-Olkin (KMO):** Measures sampling adequacy (Value $> 0.6$ is generally acceptable).
* **Bartlett’s Test of Sphericity:** Tests $H_0: \Sigma = I$. You want to **reject** this to prove variables are correlated.

### Step B: Extraction Methods

* **Principal Axis Factoring (PAF):** Analyzes shared variance; preferred when normality is violated.
* **Maximum Likelihood (ML):** Best if data is multivariate normal; allows for goodness-of-fit tests.

### Step C: Determining the Number of Factors ($m$)

1. **Kaiser Criterion:** Keep factors with eigenvalues $> 1$.
2. **Scree Plot:** Look for the "elbow" where the eigenvalues level off.
3. **Parallel Analysis:** Compare your eigenvalues to those from a random dataset of the same size.

---

## 4. Factor Rotation

The initial solution $\Lambda$ is not unique (the "Rotation Problem"). We rotate $\Lambda$ to achieve **Simple Structure** (where each variable loads highly on only one factor).
 **Varimax (Orthogonal)** Minimizes number of factors. | Focuses on explaining one general factor.
 **Promax / Oblimin (Oblique)** Factors **can be correlated**. | More realistic for human/social behaviors.
## 5. PCA vs. Factor Analysis
This is a frequent exam question for mathematicians.
| Feature | PCA | Factor Analysis (FA) |
| **Goal** | Variance Maximization. | Correlation Explanation. |
| **Model** | $PC = w_1X_1 + w_2X_2 \dots$ | $X = \lambda_1F_1 + \lambda_2F_2 + \epsilon$ |
| **Variance** | Analyzes **Total** variance. | Analyzes **Shared** (Common) variance. |
| **Latent Variable** | No; components are just math constructs. | Yes; assumes factors "cause" the variables. |

## Quick Exam Tips

* **Indeterminacy:** Unlike PCA, factor scores in FA are estimated, not calculated exactly.
* **Loading Interpretation:** Loadings are essentially the correlation between the variable and the factor (for orthogonal rotations). 
$\lambda_{ij}^2$ is the percent variance explained.
* **Heywood Case:** If an estimation results in a communality $h^2 > 1$ or negative uniqueness, it's a mathematical impossibility usually caused by too few samples or too many factors.



PCA
---

## 1. The Optimization Objective

Given a centered data matrix $X$ ($n \times p$) where each column has mean zero, we seek a direction $w \in \mathbb{R}^p$ ($||w|| = 1$) that maximizes the variance of the projected data.

The variance of the projection $Xw$ is:


$$\text{Var}(Xw) = \frac{1}{n-1} (Xw)^T(Xw) = w^T \left( \frac{X^T X}{n-1} \right) w = w^T \Sigma w$$


where $\Sigma$ is the $p \times p$ sample covariance matrix.

Using the **Rayleigh Quotient**, the maximum of $w^T \Sigma w$ subject to $||w|| = 1$ is achieved when $w$ is the eigenvector corresponding to the largest eigenvalue of $\Sigma$.

---

## 2. Eigendecomposition of $\Sigma$

Since $\Sigma$ is a symmetric, positive semi-definite matrix, the **Spectral Theorem** guarantees:


$$\Sigma = V \Lambda V^T$$

* **$V$:** An orthogonal matrix ($V^T V = I$) whose columns $v_1, v_2, \dots, v_p$ are the eigenvectors (Principal Components).
* **$\Lambda$:** A diagonal matrix of eigenvalues $\lambda_1 \ge \lambda_2 \ge \dots \ge \lambda_p \ge 0$.

### Key Mathematical Properties:

* **Total Variance:** $\text{Tr}(\Sigma) = \sum_{i=1}^p \text{Var}(X_i) = \sum_{i=1}^p \lambda_i$.
* **Proportion of Variance Explained (PVE):** For the $k$-th component, $PVE_k = \frac{\lambda_k}{\sum \lambda_i}$.
* **Orthogonality:** The principal components are orthogonal by construction, meaning the scores (projected data) are uncorrelated: $\text{Cov}(Xv_i, Xv_j) = 0$ for $i \neq j$.

---

## 3. Connection to Singular Value Decomposition (SVD)

SVD is the more numerically stable way to compute PCA without explicitly forming $\Sigma$.
For a centered matrix $X$:


$$X = U D V^T$$

* **$V$:** The right singular vectors of $X$ (identical to the eigenvectors of $\Sigma$).
* **$U$:** The left singular vectors of $X$ (eigenvectors of $XX^T$).
* **$D$:** A diagonal matrix of singular values $s_i$.

**The Identity:** The eigenvalues of $\Sigma$ are related to the singular values of $X$ by:


$$\lambda_i = \frac{s_i^2}{n-1}$$

---

## 4. PCA as Best Low-Rank Approximation

The **Eckart-Young-Mirsky Theorem** states that if we want to find a matrix $\hat{X}$ of rank $k < p$ that minimizes the Frobenius norm $||X - \hat{X}||_F$, the solution is:


$$\hat{X}_k = \sum_{i=1}^k s_i u_i v_i^T$$


This proves that PCA provides the optimal $k$-dimensional linear subspace for representing the data in terms of minimizing reconstruction error (Mean Squared Error).

---

## 5. Geometric Interpretation

PCA can be viewed as a rigid rotation of the coordinate axes to a new system where:

1. The first axis ($v_1$) lies along the direction of maximum sample variance.
2. Each subsequent axis ($v_k$) is orthogonal to the previous ones and captures the maximum remaining variance.
3. The coordinates of the data points in this new system are called **Principal Component Scores**, $Z = XV$.

---

## 6. Mathematical Constraints and Nuances

* **Mean Centering:** Crucial. If $X$ is not centered, the first PC will likely point toward the mean of the data rather than the direction of maximum variance.
* **Scaling (Correlation vs. Covariance):** PCA on the covariance matrix $\Sigma$ is sensitive to units. PCA on the **Correlation Matrix** $R$ (where $R_{ij} = \frac{\sigma_{ij}}{\sigma_i \sigma_j}$) is equivalent to performing PCA on standardized data ($z$-scores).
* **Invariance:** PCA is rotationally invariant but **not** scale-invariant.

---

## Quick Exam Identifiers

* **The "Trace" Property:** $\sum \text{Var}(X_{original}) = \sum \text{Var}(PC_{scores})$.
* **Determinant:** $|\Sigma| = \prod \lambda_i$. If any $\lambda_i = 0$, the data lives in a lower-dimensional degenerate subspace (multi-collinearity).
* **Projection Matrix:** The matrix $P_k = V_k V_k^T$ is the orthogonal projector onto the $k$-dimensional principal subspace.

Would you like me to derive the **Lagrangian** for the PCA optimization problem, or perhaps explain the **Kernel PCA** extension for non-linear dimensionality reduction?
